\subsection{The exact definite-unitary source of the positive theta factor}
\label{sec:definite-theta-factor}

\paragraph{P-F51 (the $E_6$ theta vector from a specified compact source).}
Let $K=E_6$ and $\Theta_K=\sum_{\mu\in K^\vee/K}
\theta_{K+\mu}e_\mu$ be the positive theta vector in P-F46.
There is a positive definite Hermitian space $(V,h)$ of
dimension three over $E=\mathbb Q(\zeta)$, where
$\zeta=e^{2\pi i/3}$, whose underlying trace lattice is
exactly $K$. Put $H=U(V,h)$, so that $H(\mathbb R)=U(3)$.
There are an explicitly specified compact open subgroup
$K_f\subset H(\mathbb A_f)$, a nonnegative compactly
supported automorphic function $f_K$ on
$H(\mathbb Q)\backslash H(\mathbb A)$, and three
Schwartz functions $\Phi_\mu$ such that
\[
 \theta_{K+\mu}(\tau)
   =y^{-3/2}\int_{H(\mathbb Q)\backslash H(\mathbb A)}
       \theta_{\Phi_\mu}(g_\tau,h) f_K(h)\,dh,
 \qquad \tau=x+iy.
 \tag{P-F51a}
\]
The integral is ordinary and absolutely convergent.
The normalization uses the actual measure of a specified
double class, not a genus average or an unstated mass factor.
Its right-hand group is the derived group
$SU(J)=SL_2$ of a split unitary group of signature $(1,1)$.
The Weil action in (P-F51a) is fixed by the canonical
orthogonal Schr\"odinger action; the statement does not
silently choose a central splitting for the full $U(J)$.

The underlying symplectic group contains the exact
seesaw configuration
\[
 \bigl(O(V_{\mathbb Q},B),SL_2\bigr),\qquad
 \bigl(H,U(J)\bigr),\qquad
 H\subset O(V_{\mathbb Q},B),\quad SL_2\subset U(J).
 \tag{P-F51b}
\]
Here the displayed pairs are reductive dual pairs.
The classical theta vector has weight three and the
same finite Weil type as in P-F46; it is modular on
$\Gamma(3)$ component by component. These statements
identify an actual definite-unitary source for the
positive factor, without identifying the additional
operation $\Theta_K\mapsto\delta^{-1}\Theta_K$ or its
multiplicative Borcherds lift with the specific iterated
lifts in the supplied \emph{Remark}.

\paragraph{The Hermitian lattice and the two dual pairs.}
Use the model proved in P-F46:
\[
 A_2=\{(a,b,c)\in\mathbb Z^3:a+b+c=0\},\quad
 v=(2,-1,-1)/3,\qquad
 K=A_2^3+\mathbb Z(v,v,v).
\]
On every $A_2$ factor cyclically permute the three
coordinates, and denote the resulting simultaneous
operator by $\sigma$. It preserves the Euclidean form
$B$, has $\sigma^2+\sigma+1=0$, and satisfies
$\sigma v-v\in A_2$. Thus it preserves the actual glue
lattice $K$, not just $A_2^3$. Let $\zeta$ act by
$\sigma$ on $V=K\otimes\mathbb Q$, and set
\[
 h(u,w)=\frac{1-\zeta^2}{3}B(u,w)
       +\frac{\zeta-\zeta^2}{3}B(u,\sigma w).
 \tag{P-F51c}
\]
As in the direct calculation in (P-F46l), the identities
$B(\sigma u,w)=-B(u,w)-B(u,\sigma w)$ and
$B(\sigma u,\sigma w)=B(u,w)$ give sesquilinearity,
$h(w,u)=\overline{h(u,w)}$, and
\[
       \operatorname{Tr}_{E/\mathbb Q}h=B,
       \qquad h(u,u)=B(u,u)/2.
\]
Consequently $h$ is positive definite, has rank three,
and is even on $K$. Trace duality, applied also to
$au$ for every $a\in\mathbb Z[\zeta]$, shows
$h(K,K)\subset\mathfrak D_E^{-1}$ and that its
inverse-different Hermitian dual
$\{x\in V:h(x,K)\subset\mathfrak D_E^{-1}\}$ is exactly
$K^\vee$. In particular the
quadratic exponents and discriminant cosets are unchanged.

Let $J=E^2$ have the skew-Hermitian form
\[
 j((a,b),(c,d))=a\overline d-b\overline c.
\]
On $\mathcal W=\operatorname{Res}_{E/\mathbb Q}(V\otimes_EJ)$
use the alternating form
$\Omega=\operatorname{Tr}_{E/\mathbb Q}(h\otimes j)$.
Identifying $\mathcal W$ with $V_{\mathbb Q}\oplus V_{\mathbb Q}$
gives exactly
\[
 \Omega((u_1,u_2),(w_1,w_2))
            =B(u_1,w_2)-B(u_2,w_1).
 \tag{P-F51d}
\]
This specifies the trace normalization; no factor of two
is suppressed in the passage to the oscillator representation.
Multiplication of $j$ by $i$ at the real place has one
positive and one negative eigenvalue, so its unitary
group has signature $(1,1)$.

If $A\in SL_2(E)$, then $A^tJA=J$ for the alternating
matrix $J$. The unitary condition is $A^tJ\overline A=J$.
Together these conditions imply $\overline A=A$.
Conversely every $A\in SL_2(\mathbb Q)$ satisfies both.
The same calculation over every $\mathbb Q$-algebra
proves $SU(J)=SL_2$ as algebraic groups, not only on
real points. Its action on the two summands of (P-F51d)
is the usual symplectic action. The inclusions in
(P-F51b) therefore preserve the very same symplectic form.

For completeness, the centralizers can also be checked
directly. After scalar extension to $\mathbb C$, the
unitary group $H$ is $GL_3$, and $V_{\mathbb Q}\otimes\mathbb C$
is the sum of its standard representation and its dual.
The two are nonisomorphic, and each occurs twice in
$\mathcal W\otimes\mathbb C$. Their commuting
endomorphisms are two matrices of size two; preservation
of $\Omega$ makes them inverse transposes. Descending
back to $\mathbb Q$ gives exactly $U(J)$.
The reverse centralizer is $H$ by the same calculation
with dimensions two and three exchanged. Likewise the
standard representation of $O_6$ is absolutely irreducible;
its multiplicity space in $\mathcal W$ has dimension two,
and preservation of its alternating form leaves exactly
$SL_2$. Conversely the centralizer of this $SL_2$ is
$O(V_{\mathbb Q},B)$. This proves both dual-pair assertions
and the stated seesaw configuration.

\paragraph{The splitting and the exact finite level.}
Take the standard additive character of
$\mathbb Q\backslash\mathbb A$, whose real component is
$\psi_\infty(t)=e^{2\pi it}$, and the usual self-dual
measures. The Schr\"odinger model on
$\mathcal S(V_{\mathbb Q}(\mathbb A))$ has the canonical
action
\[
        (\omega(h)\Phi)(u)=\Phi(h^{-1}u),
                   \qquad h\in O(V_{\mathbb Q},B)(\mathbb A).
 \tag{P-F51e}
\]
For the $SL_2$ unipotent $n(b)$ it has action
$\omega(n(b))\Phi(u)=\psi(bB(u,u)/2)\Phi(u)$.
Since the quadratic dimension is six, the symplectic
Weil representation descends from the double cover
to $SL_2$. These actions and the Fourier-transform
operator commute with (P-F51e), because $h$ preserves
$B$ and its self-dual measure. This fixes the action
on $H\times SL_2$ used here. At finite places these are
Kudla, \emph{Notes on the local theta correspondence},
Definition~II.4.2 and Proposition~II.4.3, pp.~37--38,
and the orthogonal splitting on p.~39. The real action
is the corresponding Gaussian Schr\"odinger action,
whose required value is computed below.

The complete finite unitary level is
\[
 K_f=\{h\in H(\mathbb A_f):h\widehat K=\widehat K,
                  \ h\text{ fixes }\widehat K^\vee/\widehat K
                          \text{ pointwise}\}.
 \tag{P-F51f}
\]
This is compact open: it is the kernel of a finite
action of the lattice stabilizer. Away from three,
$K_p=K\otimes\mathbb Z_p$ is self-dual, so the discriminant condition
adds no restriction there. For each class $\mu$ put
\[
 \Phi_\mu(u)=e^{-\pi B(u_\infty,u_\infty)}
                 \mathbf1_{\mu+\widehat K}(u_f).
 \tag{P-F51g}
\]
The Gaussian is $H(\mathbb R)$-invariant and the
finite indicator is $K_f$-invariant. Requiring each
coset to be fixed, rather than merely the lattice,
is what makes this simultaneous assertion immediate.

In the labels $\mu\in\mathbb Z/3\mathbb Z$ with
$q(\mu)=2\mu^2/3$, the classical Weil matrices are
\[
 \rho_K(T)e_\mu=e^{2\pi i(2\mu^2/3)}e_\mu,
 \qquad
 \rho_K(S)e_\mu=\frac{i}{\sqrt3}
                  \sum_\nu e^{-2\pi i\mu\nu/3}e_\nu.
 \tag{P-F51h}
\]
Indeed the first formula follows term by term; the
second is Poisson summation, with the factor
$(-i\tau)^3=i\tau^3$, followed if necessary by
$\nu\mapsto-\nu$, since opposite cosets have equal
theta series. These are the same convention as P-F46.
The finite quadratic module has level three.
Scheithauer, Proposition~4.3, printed p.~14, gives
trivial action of $\Gamma(3)$. His displayed convention
is dual to this one; taking a dual leaves this kernel
assertion unchanged. Thus each component is a genuine
weight-three modular form on $\Gamma(3)$, without
claiming an unspecified scalar level for the full
unitary group.

A full $U(J)$ formulation would additionally select
unitary splitting characters with restrictions
$\chi_V|_{\mathbb A_\mathbb Q^\times}=\varepsilon_{E/\mathbb Q}^3$
and $\chi_J|_{\mathbb A_\mathbb Q^\times}=\varepsilon_{E/\mathbb Q}^2$.
These are the conditions in Harris--Li--Sun,
equation~(1.1.1), p.~3, and Section~1.2, p.~4.
No such choice, nor compatibility with a different
iterated lift's choices, is silently supplied by
(P-F51a); the derived-group and orthogonal splitting
above already suffice for its exact identity.

\paragraph{Selecting the individual lattice class.}
Let $C_K$ be the image of
$H(\mathbb R)K_f$ in $[H]=H(\mathbb Q)\backslash H(\mathbb A)$.
It is a nonempty compact open subset: its preimage is
$H(\mathbb Q)H(\mathbb R)K_f$, which is open, and
$H(\mathbb R)K_f$ is compact. For any fixed Haar
quotient measure its volume $c_K$ is finite and positive.
Set
\[
                   f_K=c_K^{-1}\mathbf1_{C_K}.
 \tag{P-F51i}
\]
It is an automorphic function, is invariant under
$H(\mathbb R)K_f$ on the right, and has compact support.
In particular it is square-integrable. There is no
assertion that this lattice-class selector lies in
one preselected irreducible automorphic representation.

Define the theta kernel in the stated Weil model by
\[
 \theta_{\Phi_\mu}(g,h)=
       \sum_{u\in V_{\mathbb Q}(\mathbb Q)}
                         (\omega(g)\omega(h)\Phi_\mu)(u).
 \tag{P-F51j}
\]
The automorphic theta sum and its integration convention
are the constructions of Harris--Li--Sun,
equations~(1.2.4)--(1.2.5), p.~4, applied to the
explicitly fixed commuting actions above.
Rational $H$-invariance also follows
directly by changing the summation variable.
For $h\in C_K$, write a representative as
$\gamma h_\infty k$ with $\gamma\in H(\mathbb Q)$,
$h_\infty\in H(\mathbb R)$, and $k\in K_f$.
Rational invariance and (P-F51g) give
$\theta_{\Phi_\mu}(g,h)=\theta_{\Phi_\mu}(g,1)$.
Thus its integral against (P-F51i) is exactly
$\theta_{\Phi_\mu}(g,1)$. This proves the normalization
and shows why a genus average has not been substituted
for the requested individual lattice theta series.

Finally take $g_\tau=n(x)m(\sqrt y)$ at infinity and
identity at all finite places. In real quadratic
dimension six the dilation factor is
$(\sqrt y)^{6/2}=y^{3/2}$. The two real operators give
\[
 (\omega(g_\tau)\Phi_\mu)(u)
     =y^{3/2}e^{2\pi i\tau B(u_\infty,u_\infty)/2}
                            \mathbf1_{\mu+\widehat K}(u_f).
\]
The finite condition on a rational vector is exactly
$u\in K+\mu$; this follows coordinate by coordinate in
an integral basis of $K$. Summing proves (P-F51a),
with the exact power of $y$ and the exact exponent.
The Gaussian absolute-value series is convergent,
and on the compact support $C_K$ the kernel is constant;
therefore the integral used in the identity is absolutely
convergent as well.\hfill$\square$

\paragraph{The precise overlap with the supplied theta method.}
This is a source-map statement for the actual factor
$\Theta_{E_6}$ in (P-F46a): it comes from a specified
nonnegative function on a definite $U(3)$ quotient,
with coset level and measure normalization fixed.
The compatible seesaw in (P-F51b) is also explicit.
Its right unitary space is split of signature $(1,1)$;
it is not a pair of totally definite spaces of the
kind named in item~8 of the supplied \emph{Remark}.
This identifies the exact pair involved rather than
asserting that signature differences prohibit all
comparisons. A further seesaw or representation-theoretic
identity would have to match the selected vectors and
splittings to the particular $U(X_j)$ in that note.

Two additional operations remain after this positive
definite factor: multiplication by $\delta^{-1}$ and
the regularized multiplicative lift identified in
P-F46. The former introduces the displayed negative
Fourier powers; the latter outputs the single cubic
discriminant, rather than every invariant-ring graded
piece. The ordinary convergence just proved for
$\Theta_{E_6}$ does not remove the explicit divergence
proved in P-F46 for $g=\Theta_{E_6}/\delta$.
No simultaneous-stretch Hilbert identity, strict-index
bound, or comparison with the supplied iterated lift's
Bessel period is inferred from (P-F51a).

Primary source locators: Stephen Kudla,
\emph{Notes on the local theta correspondence},
Definition~II.4.2 and Proposition~II.4.3, pp.~37--38;
orthogonal splitting, p.~39; seesaw definition,
Section~IV.1, p.~59,
\url{https://www.math.toronto.edu/~skudla/castle.pdf}.
Michael Harris, Jian-Shu Li and Binyong Sun,
\emph{Theta correspondences for close unitary groups},
equation~(1.1.1), p.~3, and equations~(1.2.1)--(1.2.5),
p.~4,
\url{https://webusers.imj-prg.fr/~michael.harris/closeunitaryweb.pdf}.
Nils Scheithauer, \emph{The Weil representation of $SL_2(\mathbb Z)$
and some applications}, Proposition~4.3, p.~14,
\url{https://www.mathematik.tu-darmstadt.de/media/algebra/homepages/scheithauer/publications/Weil_representation.pdf}.
